\documentclass[11pt]{article}
\usepackage{amsmath,amssymb,geometry,enumitem,longtable,array,tcolorbox,xcolor,hyperref}
\geometry{margin=0.75in}
\setlength{\parindent}{0pt}
\setlength{\parskip}{0.55em}

\begin{document}

\title{CAT Quant Verified Practice Batch 01}
\author{Original PYQ-Inspired Practice Set}
\date{}
\maketitle

\section*{1. Instructions}

\begin{itemize}
  \item This is an original practice batch inspired by CAT Quant PYQ patterns. It is not an official, leaked, or guaranteed CAT paper.
  \item Total questions: 15 (9 MCQ + 6 TITA).
  \item Suggested time: 30 to 40 minutes.
  \item Calculators are not assumed.
  \item For MCQs, exactly one option is correct.
  \item For TITA questions, write an exact numeric or integer answer (no approximation unless explicitly defined).
\end{itemize}

\section*{2. Question Paper}

\subsection*{Question 1}
\textbf{Topic:} Arithmetic \quad
\textbf{Subtopic:} Profit, Loss \& Discount \quad
\textbf{Difficulty:} Medium-Hard \quad
\textbf{Type:} MCQ \quad
\textbf{Estimated Time:} 3 min \quad
\textbf{PYQ-Inspired Archetype:} Marked-price chain with hidden free-item adjustment.

\begin{tcolorbox}
A trader marks his goods 50\% above cost price. On the marked price he offers a discount of \(d\%\). In addition, to clear stock he runs a ``Buy 5, Get 1 Free'' scheme. After all of this, his net profit on the entire transaction is exactly 20\%. The value of \(d\) is:
\end{tcolorbox}

\begin{enumerate}[label=(\Alph*)]
\item 4
\item 5
\item 6
\item 8
\end{enumerate}

\subsection*{Question 2}
\textbf{Topic:} Arithmetic \quad
\textbf{Subtopic:} Time, Speed \& Distance \quad
\textbf{Difficulty:} Hard \quad
\textbf{Type:} MCQ \quad
\textbf{Estimated Time:} 4 min \quad
\textbf{PYQ-Inspired Archetype:} Two-body meet with post-meeting time ratio.

\begin{tcolorbox}
Two trains \(A\) and \(B\) leave stations \(P\) and \(Q\) simultaneously and travel towards each other on the same track. After they meet, \(A\) takes 9 hours to reach \(Q\), while \(B\) takes 16 hours to reach \(P\). If the speed of \(A\) is 60 km/h, the distance \(PQ\) (in km) is:
\end{tcolorbox}

\begin{enumerate}[label=(\Alph*)]
\item 1050
\item 1260
\item 1470
\item 840
\end{enumerate}

\subsection*{Question 3}
\textbf{Topic:} Algebra \quad
\textbf{Subtopic:} Logarithms \quad
\textbf{Difficulty:} Medium \quad
\textbf{Type:} MCQ \quad
\textbf{Estimated Time:} 3 min \quad
\textbf{PYQ-Inspired Archetype:} Logarithms with different bases of the same prime.

\begin{tcolorbox}
If \(x>0\) and \(\log_{2}x + \log_{4}x + \log_{8}x = 11\), then \(x\) equals:
\end{tcolorbox}

\begin{enumerate}[label=(\Alph*)]
\item 32
\item 64
\item 128
\item 256
\end{enumerate}

\subsection*{Question 4}
\textbf{Topic:} Number System \quad
\textbf{Subtopic:} Modular Arithmetic / Remainders \quad
\textbf{Difficulty:} Hard \quad
\textbf{Type:} TITA \quad
\textbf{Estimated Time:} 4 min \quad
\textbf{PYQ-Inspired Archetype:} Large-power remainder via Fermat's little theorem.

\begin{tcolorbox}
Find the remainder when \(7^{100} + 11^{100}\) is divided by 13.
\end{tcolorbox}

\textbf{Type:} TITA

\subsection*{Question 5}
\textbf{Topic:} Geometry \quad
\textbf{Subtopic:} Circles \& Tangents \quad
\textbf{Difficulty:} Medium-Hard \quad
\textbf{Type:} MCQ \quad
\textbf{Estimated Time:} 3 min \quad
\textbf{PYQ-Inspired Archetype:} Direct common tangent of two externally tangent circles.

\begin{tcolorbox}
Two circles of radii 9 cm and 16 cm touch each other externally at a point \(T\). A direct (external) common tangent touches the larger circle at \(P\) and the smaller circle at \(Q\). The length of \(PQ\) (in cm) is:
\end{tcolorbox}

\begin{enumerate}[label=(\Alph*)]
\item 20
\item 24
\item 25
\item 12
\end{enumerate}

\subsection*{Question 6}
\textbf{Topic:} Arithmetic \quad
\textbf{Subtopic:} Mixtures \& Replacement \quad
\textbf{Difficulty:} Medium \quad
\textbf{Type:} TITA \quad
\textbf{Estimated Time:} 3 min \quad
\textbf{PYQ-Inspired Archetype:} Repeated replacement of liquid (geometric depletion).

\begin{tcolorbox}
A vessel contains 80 litres of pure milk. From it, 20 litres of the liquid is drawn out and replaced with water; this process is repeated two more times (three times in all). The quantity of pure milk (in litres) remaining in the vessel at the end is:
\end{tcolorbox}

\textbf{Type:} TITA

\subsection*{Question 7}
\textbf{Topic:} Algebra \quad
\textbf{Subtopic:} Quadratic Equations / Integer Roots \quad
\textbf{Difficulty:} Hard \quad
\textbf{Type:} TITA \quad
\textbf{Estimated Time:} 4 min \quad
\textbf{PYQ-Inspired Archetype:} Integer-root parameter counting via factor decomposition.

\begin{tcolorbox}
For how many distinct integer values of \(k\) does the equation
\[
x^{2} - kx + (k+8) = 0
\]
have both roots as integers?
\end{tcolorbox}

\textbf{Type:} TITA

\subsection*{Question 8}
\textbf{Topic:} Number System \quad
\textbf{Subtopic:} Divisibility \& Inclusion–Exclusion \quad
\textbf{Difficulty:} Hard \quad
\textbf{Type:} MCQ \quad
\textbf{Estimated Time:} 4 min \quad
\textbf{PYQ-Inspired Archetype:} ``Exactly \(k\) of these primes'' counting trap.

\begin{tcolorbox}
The number of positive integers \(n \le 1000\) that are divisible by exactly three of the primes 2, 3, 5 and 7 is:
\end{tcolorbox}

\begin{enumerate}[label=(\Alph*)]
\item 59
\item 63
\item 67
\item 71
\end{enumerate}

\subsection*{Question 9}
\textbf{Topic:} Algebra \quad
\textbf{Subtopic:} Functional Equations \quad
\textbf{Difficulty:} Medium-Hard \quad
\textbf{Type:} MCQ \quad
\textbf{Estimated Time:} 3 min \quad
\textbf{PYQ-Inspired Archetype:} Substitution-pair functional equation.

\begin{tcolorbox}
A function \(f:\mathbb{R}\to\mathbb{R}\) satisfies
\[
f(x) + 2f(1-x) = 3x^{2} \quad \text{for every real } x.
\]
The value of \(f(2)\) is:
\end{tcolorbox}

\begin{enumerate}[label=(\Alph*)]
\item \(-2\)
\item 0
\item 2
\item 4
\end{enumerate}

\subsection*{Question 10}
\textbf{Topic:} Arithmetic \quad
\textbf{Subtopic:} Boats \& Streams \quad
\textbf{Difficulty:} Hard \quad
\textbf{Type:} TITA \quad
\textbf{Estimated Time:} 5 min \quad
\textbf{PYQ-Inspired Archetype:} Stream-speed perturbation with two scenarios.

\begin{tcolorbox}
A boat takes 6 hours to complete a round trip of 36 km downstream and back upstream. If the speed of the stream is doubled (boat speed in still water unchanged), the same round trip takes 8 hours. The speed of the boat in still water (in km/h) is:
\end{tcolorbox}

\textbf{Type:} TITA

\subsection*{Question 11}
\textbf{Topic:} Geometry \quad
\textbf{Subtopic:} Coordinate Geometry / Regions \quad
\textbf{Difficulty:} Medium-Hard \quad
\textbf{Type:} MCQ \quad
\textbf{Estimated Time:} 3 min \quad
\textbf{PYQ-Inspired Archetype:} Region between a rhombus and a circle.

\begin{tcolorbox}
Find the area of the region in the \(xy\)-plane defined by
\[
\{(x,y) : |x| + |y| \le 4 \text{ and } x^{2} + y^{2} \ge 4\}.
\]
\end{tcolorbox}

\begin{enumerate}[label=(\Alph*)]
\item \(32 - 4\pi\)
\item \(16 - 4\pi\)
\item \(32 - 2\pi\)
\item \(16 - 2\pi\)
\end{enumerate}

\subsection*{Question 12}
\textbf{Topic:} Arithmetic \quad
\textbf{Subtopic:} Averages \& Extremes with Median Constraint \quad
\textbf{Difficulty:} Medium-Hard \quad
\textbf{Type:} TITA \quad
\textbf{Estimated Time:} 3 min \quad
\textbf{PYQ-Inspired Archetype:} Maximising one term under sum, distinctness and median constraints.

\begin{tcolorbox}
Seven distinct positive integers have an average of 12 and a median of 10. The maximum possible value of the largest of these seven integers is:
\end{tcolorbox}

\textbf{Type:} TITA

\subsection*{Question 13}
\textbf{Topic:} Arithmetic \quad
\textbf{Subtopic:} Time \& Work (Alternate Days) \quad
\textbf{Difficulty:} Medium-Hard \quad
\textbf{Type:} TITA \quad
\textbf{Estimated Time:} 4 min \quad
\textbf{PYQ-Inspired Archetype:} Alternate-day work with fractional finish.

\begin{tcolorbox}
\(A\) alone can finish a job in 20 days and \(B\) alone can finish it in 25 days. They work on alternate days, with \(A\) working on Day 1, \(B\) on Day 2, \(A\) on Day 3, and so on. The total time (in days) taken to complete the job is:
\end{tcolorbox}

\textbf{Type:} TITA

\subsection*{Question 14}
\textbf{Topic:} Algebra \quad
\textbf{Subtopic:} Logarithmic Inequalities \quad
\textbf{Difficulty:} Medium-Hard \quad
\textbf{Type:} MCQ \quad
\textbf{Estimated Time:} 3 min \quad
\textbf{PYQ-Inspired Archetype:} Log-inequality with domain trap.

\begin{tcolorbox}
The number of integer values of \(x\) satisfying
\[
\log_{2}\!\big(x^{2} - 5x + 6\big) \le 1
\]
is:
\end{tcolorbox}

\begin{enumerate}[label=(\Alph*)]
\item 1
\item 2
\item 3
\item 4
\end{enumerate}

\subsection*{Question 15}
\textbf{Topic:} Algebra \quad
\textbf{Subtopic:} Sequences \& Series \quad
\textbf{Difficulty:} Medium \quad
\textbf{Type:} TITA \quad
\textbf{Estimated Time:} 3 min \quad
\textbf{PYQ-Inspired Archetype:} Sum of cumulative-sum sequence (triangular numbers).

\begin{tcolorbox}
Let \(T_{n} = 1 + 2 + 3 + \cdots + n\). Find the value of
\[
T_{1} + T_{2} + T_{3} + \cdots + T_{20}.
\]
\end{tcolorbox}

\textbf{Type:} TITA

\newpage
\section*{3. Answer Key}

\begin{longtable}{|c|l|c|c|c|}
\hline
\textbf{Q.No.} & \textbf{Topic} & \textbf{Type} & \textbf{Answer} & \textbf{Difficulty} \\
\hline
1 & Arithmetic – P/L \& Discount & MCQ & (A) 4 & Medium-Hard \\ \hline
2 & Arithmetic – TSD & MCQ & (B) 1260 & Hard \\ \hline
3 & Algebra – Logarithms & MCQ & (B) 64 & Medium \\ \hline
4 & Number System – Remainders & TITA & 12 & Hard \\ \hline
5 & Geometry – Tangents & MCQ & (B) 24 & Medium-Hard \\ \hline
6 & Arithmetic – Mixtures & TITA & 33.75 & Medium \\ \hline
7 & Algebra – Quadratics & TITA & 4 & Hard \\ \hline
8 & Number System – Divisibility & MCQ & (B) 63 & Hard \\ \hline
9 & Algebra – Functions & MCQ & (A) \(-2\) & Medium-Hard \\ \hline
10 & Arithmetic – Boats \& Streams & TITA & 13 & Hard \\ \hline
11 & Geometry – Region/Area & MCQ & (A) \(32-4\pi\) & Medium-Hard \\ \hline
12 & Arithmetic – Averages & TITA & 45 & Medium-Hard \\ \hline
13 & Arithmetic – Time \& Work & TITA & 22.2 (= 111/5) & Medium-Hard \\ \hline
14 & Algebra – Log Inequality & MCQ & (B) 2 & Medium-Hard \\ \hline
15 & Algebra – Series & TITA & 1540 & Medium \\ \hline
\end{longtable}

\section*{4. Detailed Solutions}

\subsection*{Solution 1}
\textbf{Answer:} (A) 4

\textbf{Detailed Solution:}

Step 1: Let cost price per item be \(C\). Marked price \( = 1.5C\).

Step 2: In the ``Buy 5, Get 1 Free'' scheme the customer pays for 5 items but receives 6. So for every 6 items the trader effectively sells, the revenue is the selling price of 5 items.

Step 3: Selling price per item after discount \(d\%\) is \(1.5C(1 - d/100)\). Revenue for 6 items \(= 5 \times 1.5C(1 - d/100) = 7.5C\,(1 - d/100)\).

Step 4: Cost of 6 items \(= 6C\). For a 20\% profit, revenue \(= 1.2 \times 6C = 7.2C\).

Step 5: Set \(7.5C(1-d/100) = 7.2C\). So \(1 - d/100 = 7.2/7.5 = 0.96\), giving \(d = 4\).

\textbf{Why this is CAT-level:} Mixes two distinct discount mechanisms (percentage discount on MP and a free-item scheme) — a classic CAT layering trap.

\textbf{Common Wrong Approach:} Treating ``Buy 5 Get 1 Free'' as a flat one-sixth (\(\approx 16.67\%\)) extra discount on the selling price and adding it to \(d\) linearly leads to wrong answers like 5 or 6.

\subsection*{Solution 2}
\textbf{Answer:} (B) 1260

\textbf{Detailed Solution:}

Step 1: For two bodies starting from opposite ends and meeting somewhere in between, if they take \(t_A\) and \(t_B\) hours to reach the opposite station \emph{after} the meeting, then \(\dfrac{\text{speed}_A}{\text{speed}_B} = \sqrt{\dfrac{t_B}{t_A}}\).

Step 2: \(\dfrac{\text{speed}_A}{\text{speed}_B} = \sqrt{\dfrac{16}{9}} = \dfrac{4}{3}\). With speed of \(A = 60\), speed of \(B = 45\) km/h.

Step 3: Let \(t\) be the time before they meet. Distance covered by \(B\) before meeting \(= 45t\). After meeting, \(A\) covers the same \(45t\) in 9 hours at 60 km/h: \(45t = 60 \times 9 \Rightarrow t = 12\).

Step 4: Distance \(PQ = 60 \times 12 + 45 \times 12 = 720 + 540 = 1260\) km.

Step 5: Verification: \(B\) covers \(60 \times 12 = 720\) km after meeting in 16 hours, i.e., at \(720/16 = 45\) km/h. Consistent.

\textbf{Why this is CAT-level:} Requires the non-obvious square-root relation between pre- and post-meeting times.

\textbf{Common Wrong Approach:} Setting up speed ratio as \(16:9\) directly (without square root) gives 1470 km — a classic CAT trap option.

\subsection*{Solution 3}
\textbf{Answer:} (B) 64

\textbf{Detailed Solution:}

Step 1: Convert each log to base 2: \(\log_{4}x = \frac{1}{2}\log_{2}x,\ \log_{8}x = \frac{1}{3}\log_{2}x\).

Step 2: The equation becomes \(\log_{2}x\Big(1 + \tfrac{1}{2} + \tfrac{1}{3}\Big) = 11\), i.e., \(\log_{2}x \cdot \tfrac{11}{6} = 11\).

Step 3: So \(\log_{2}x = 6\) and \(x = 2^{6} = 64\).

\textbf{Why this is CAT-level:} Standard log manipulation, but the coefficient \(11/6\) collapses cleanly only if the base conversion is done correctly.

\textbf{Common Wrong Approach:} Adding the bases (\(2 + 4 + 8 = 14\)) and writing \(\log_{14}x = 11\) — a trap many students fall into.

\subsection*{Solution 4}
\textbf{Answer:} 12

\textbf{Detailed Solution:}

Step 1: By Fermat's little theorem, for any \(a\) coprime to 13, \(a^{12} \equiv 1 \pmod{13}\). Since \(100 = 12 \cdot 8 + 4\), \(a^{100} \equiv a^{4} \pmod{13}\).

Step 2: Compute \(7^{4} \pmod{13}\). \(7^{2} = 49 \equiv 10\); \(7^{4} \equiv 10^{2} = 100 \equiv 9 \pmod{13}\).

Step 3: Compute \(11^{4} \pmod{13}\). Note \(11 \equiv -2 \pmod{13}\), so \(11^{4} \equiv (-2)^{4} = 16 \equiv 3 \pmod{13}\).

Step 4: Sum: \(7^{100} + 11^{100} \equiv 9 + 3 = 12 \pmod{13}\).

Step 5: The remainder is \(\boxed{12}\).

\textbf{Why this is CAT-level:} Requires recognising Fermat's exponent reduction and switching to a negative residue to simplify.

\textbf{Common Wrong Approach:} Reducing the exponent modulo 13 instead of 12 (Fermat's order), giving the wrong intermediate residue.

\subsection*{Solution 5}
\textbf{Answer:} (B) 24

\textbf{Detailed Solution:}

Step 1: For two circles with centres \(O_1\) and \(O_2\) and radii \(r_1, r_2\), the length of a direct (external) common tangent is
\(L = \sqrt{d^{2} - (r_1 - r_2)^{2}}\),
where \(d = O_1O_2\).

Step 2: Since the circles touch externally, \(d = r_1 + r_2 = 9 + 16 = 25\).

Step 3: \(|r_1 - r_2| = 7\). So \(L = \sqrt{25^{2} - 7^{2}} = \sqrt{625 - 49} = \sqrt{576} = 24\).

\textbf{Why this is CAT-level:} The student must distinguish the formulae for direct vs.\ transverse common tangents.

\textbf{Common Wrong Approach:} Using \(\sqrt{d^{2} - (r_1+r_2)^{2}}\) (the transverse tangent formula). For externally tangent circles this gives 0, leading some to pick 0/12 incorrectly.

\subsection*{Solution 6}
\textbf{Answer:} 33.75 litres

\textbf{Detailed Solution:}

Step 1: For repeated replacement, the milk remaining after \(n\) operations is
\(\text{Milk}_{\text{left}} = V\left(1 - \dfrac{x}{V}\right)^{n}\),
where \(V\) is the vessel capacity, \(x\) is the volume replaced each time.

Step 2: Here \(V = 80\), \(x = 20\), \(n = 3\). Fraction left \(= (1 - 20/80)^{3} = (3/4)^{3} = 27/64\).

Step 3: Milk remaining \(= 80 \cdot 27/64 = 27 \cdot 80/64 = 27 \cdot 1.25 = 33.75\) litres.

\textbf{Why this is CAT-level:} The classical formula is well known, but the trap is mis-counting the number of replacements (the phrase ``two more times'' is deliberate).

\textbf{Common Wrong Approach:} Computing only two replacements (giving \(80 \cdot 9/16 = 45\)) — a common mis-reading.

\subsection*{Solution 7}
\textbf{Answer:} 4

\textbf{Detailed Solution:}

Step 1: Let the integer roots be \(r\) and \(s\). Then \(r+s = k\) and \(rs = k+8\).

Step 2: Eliminate \(k\): \(rs - (r+s) = 8\), i.e., \(rs - r - s + 1 = 9\), so \((r-1)(s-1) = 9\).

Step 3: List integer factor pairs of 9: \((1,9), (9,1), (3,3), (-1,-9), (-9,-1), (-3,-3)\).

Step 4: Corresponding \((r,s)\): \((2,10), (10,2), (4,4), (0,-8), (-8,0), (-2,-2)\). The values of \(k = r+s\) are \(12, 12, 8, -8, -8, -4\).

Step 5: Distinct \(k\) values are \(\{12, 8, -8, -4\}\), giving \(\boxed{4}\) integer values.

\textbf{Why this is CAT-level:} The transformation \((r-1)(s-1)=9\) is the key insight; brute searching by trying values of \(k\) is slow.

\textbf{Common Wrong Approach:} Double-counting symmetric pairs as distinct \(k\)-values, leading to answer 6.

\subsection*{Solution 8}
\textbf{Answer:} (B) 63

\textbf{Detailed Solution:}

Step 1: We need integers \(n \le 1000\) divisible by exactly three of \(\{2,3,5,7\}\). There are \(\binom{4}{3} = 4\) ways to choose which three primes divide \(n\); the fourth prime must \emph{not} divide \(n\).

Step 2: For each choice, ``divisible by these three but not the fourth'' equals (multiples of product of three) minus (multiples of product of all four = 210).

Step 3: Compute:
\begin{itemize}
  \item \(\{2,3,5\}\) (not 7): \(\lfloor 1000/30\rfloor - \lfloor 1000/210\rfloor = 33 - 4 = 29\).
  \item \(\{2,3,7\}\) (not 5): \(\lfloor 1000/42\rfloor - \lfloor 1000/210\rfloor = 23 - 4 = 19\).
  \item \(\{2,5,7\}\) (not 3): \(\lfloor 1000/70\rfloor - \lfloor 1000/210\rfloor = 14 - 4 = 10\).
  \item \(\{3,5,7\}\) (not 2): \(\lfloor 1000/105\rfloor - \lfloor 1000/210\rfloor = 9 - 4 = 5\).
\end{itemize}

Step 4: Total \(= 29 + 19 + 10 + 5 = 63\).

\textbf{Why this is CAT-level:} A clean illustration of inclusion–exclusion combined with ``exactly \(k\)'' conditions.

\textbf{Common Wrong Approach:} Forgetting to subtract multiples of 210 and reporting 79 (\(=33+23+14+9\)).

\subsection*{Solution 9}
\textbf{Answer:} (A) \(-2\)

\textbf{Detailed Solution:}

Step 1: Substitute \(x = 2\) in \(f(x) + 2f(1-x) = 3x^{2}\):
\(f(2) + 2f(-1) = 12\). \quad (i)

Step 2: Substitute \(x = -1\):
\(f(-1) + 2f(2) = 3\). \quad (ii)

Step 3: From (ii): \(f(-1) = 3 - 2f(2)\). Substitute into (i):
\(f(2) + 2(3 - 2f(2)) = 12 \Rightarrow f(2) + 6 - 4f(2) = 12 \Rightarrow -3f(2) = 6\).

Step 4: Hence \(f(2) = -2\).

\textbf{Why this is CAT-level:} Standard substitution-pair technique for functional equations — a recurring CAT pattern.

\textbf{Common Wrong Approach:} Trying to guess \(f(x) = ax^{2} + bx + c\) and getting algebra wrong on the cross terms.

\subsection*{Solution 10}
\textbf{Answer:} 13 km/h

\textbf{Detailed Solution:}

Step 1: Let \(b\) be the boat's still-water speed and \(s\) the stream speed. Time for the round trip is \(\dfrac{36}{b+s} + \dfrac{36}{b-s} = \dfrac{72b}{b^{2}-s^{2}}\).

Step 2: First scenario: \(\dfrac{72b}{b^{2}-s^{2}} = 6 \Rightarrow b^{2}-s^{2} = 12b\). \quad (i)

Step 3: Second scenario (stream doubled): \(\dfrac{72b}{b^{2}-4s^{2}} = 8 \Rightarrow b^{2}-4s^{2} = 9b\). \quad (ii)

Step 4: From (i): \(s^{2} = b^{2}-12b\). Substitute into (ii):
\(b^{2} - 4(b^{2}-12b) = 9b \Rightarrow -3b^{2} + 48b = 9b \Rightarrow 3b^{2} = 39b \Rightarrow b = 13\) (\(b>0\)).

Step 5: Check: \(s^{2} = 169 - 156 = 13\), so \(s = \sqrt{13}\) (positive and less than \(b\)).

\textbf{Why this is CAT-level:} Two non-linear equations in two unknowns reduce neatly only after isolating \(s^{2}\).

\textbf{Common Wrong Approach:} Assuming \(b/s\) ratio is preserved and writing 12 km/h.

\subsection*{Solution 11}
\textbf{Answer:} (A) \(32 - 4\pi\)

\textbf{Detailed Solution:}

Step 1: The region \(|x|+|y| \le 4\) is a square (rotated \(45^\circ\)) with vertices \((\pm 4, 0)\) and \((0, \pm 4)\). Its diagonals are both 8, so area \(= \tfrac{1}{2}\cdot 8 \cdot 8 = 32\).

Step 2: The region \(x^{2}+y^{2} \ge 4\) excludes a disc of radius 2 (area \(4\pi\)) centred at the origin.

Step 3: Check that this disc lies completely inside the rhombus: distance from origin to each side of \(|x|+|y|=4\) is \(4/\sqrt{2} = 2\sqrt{2} \approx 2.83 > 2\). So the disc is fully inside the square.

Step 4: Required area \(= 32 - 4\pi\).

\textbf{Why this is CAT-level:} Combines two standard regions and requires the containment check.

\textbf{Common Wrong Approach:} Forgetting to check containment and subtracting only the disc area inside one quadrant (giving \(32 - \pi\)).

\subsection*{Solution 12}
\textbf{Answer:} 45

\textbf{Detailed Solution:}

Step 1: Let the 7 distinct positive integers, in increasing order, be \(a_1 < a_2 < a_3 < a_4 < a_5 < a_6 < a_7\). The median is the 4th term: \(a_4 = 10\). The sum is \(7 \times 12 = 84\).

Step 2: To maximise \(a_7\), minimise all other six terms.

Step 3: The three smallest distinct positive integers smaller than 10 are 1, 2, 3 — so \(a_1, a_2, a_3 = 1, 2, 3\).

Step 4: For \(a_5\) and \(a_6\), the smallest distinct values strictly greater than 10 are 11 and 12.

Step 5: Sum of the six fixed terms: \(1+2+3+10+11+12 = 39\). Thus \(a_7 = 84 - 39 = 45\).

\textbf{Why this is CAT-level:} The student must respect strict distinctness on \emph{both} sides of the median.

\textbf{Common Wrong Approach:} Using \(a_5 = a_4 = 10\) or allowing duplicates; or forgetting positivity and picking \(a_1 = 0\).

\subsection*{Solution 13}
\textbf{Answer:} 22.2 days (i.e., \(\tfrac{111}{5}\) days)

\textbf{Detailed Solution:}

Step 1: In one day, \(A\) does \(\tfrac{1}{20}\) of the work; \(B\) does \(\tfrac{1}{25}\).

Step 2: In each two-day cycle (one day of \(A\) followed by one day of \(B\)), work done \(= \tfrac{1}{20} + \tfrac{1}{25} = \tfrac{5+4}{100} = \tfrac{9}{100}\).

Step 3: After 11 complete cycles (22 days), work done \(= 11 \times \tfrac{9}{100} = \tfrac{99}{100}\). Remaining \(= \tfrac{1}{100}\).

Step 4: Day 23 is \(A\)'s turn; \(A\) does \(\tfrac{1}{20}\) per full day. Time needed for the remaining \(\tfrac{1}{100}\) is \(\dfrac{1/100}{1/20} = \tfrac{1}{5}\) day.

Step 5: Total time \(= 22 + \tfrac{1}{5} = 22.2\) days.

\textbf{Why this is CAT-level:} The trap is to round up to ``Day 23'' as a count; the precise total uses the fractional portion of the day.

\textbf{Common Wrong Approach:} Reporting 23 days because work spills into the 23rd day.

\subsection*{Solution 14}
\textbf{Answer:} (B) 2

\textbf{Detailed Solution:}

Step 1: Domain condition: \(x^{2}-5x+6 > 0 \Rightarrow (x-2)(x-3) > 0 \Rightarrow x < 2\) or \(x > 3\).

Step 2: Inequality: \(\log_{2}(x^{2}-5x+6) \le 1 \Rightarrow x^{2}-5x+6 \le 2 \Rightarrow x^{2}-5x+4 \le 0\)
\(\Rightarrow (x-1)(x-4) \le 0 \Rightarrow 1 \le x \le 4\).

Step 3: Combine: \((-\infty,2)\cup(3,\infty)\) intersected with \([1,4]\) gives \([1,2) \cup (3,4]\).

Step 4: Integers in this set: \(x = 1\) (in \([1,2)\)) and \(x = 4\) (in \((3,4]\)). \(x=2\) and \(x=3\) are excluded by the domain. Count \(= 2\).

\textbf{Why this is CAT-level:} Both a strict domain and the inequality must be enforced; missing the open–closed endpoints is the trap.

\textbf{Common Wrong Approach:} Including \(x = 2, 3\) and answering 4.

\subsection*{Solution 15}
\textbf{Answer:} 1540

\textbf{Detailed Solution:}

Step 1: \(T_n = \dfrac{n(n+1)}{2}\).

Step 2: Sum required \(= \displaystyle\sum_{n=1}^{20} \dfrac{n(n+1)}{2} = \dfrac{1}{2}\sum_{n=1}^{20}(n^{2}+n) = \dfrac{1}{2}\!\left[\sum_{n=1}^{20} n^{2} + \sum_{n=1}^{20} n\right]\).

Step 3: \(\displaystyle\sum_{n=1}^{20} n^{2} = \dfrac{20 \cdot 21 \cdot 41}{6} = 2870\); \quad \(\displaystyle\sum_{n=1}^{20} n = \dfrac{20 \cdot 21}{2} = 210\).

Step 4: Sum \(= \dfrac{1}{2}(2870 + 210) = \dfrac{3080}{2} = 1540\).

Step 5: (Cross-check via tetrahedral formula: \(\sum_{n=1}^{N} T_n = \dfrac{N(N+1)(N+2)}{6} = \dfrac{20 \cdot 21 \cdot 22}{6} = 1540\). Matches.)

\textbf{Why this is CAT-level:} Tests knowledge of standard summations and recognition of the tetrahedral identity.

\textbf{Common Wrong Approach:} Using \(\sum T_n = \big(\sum n\big)^{2}/2\) or some incorrect closed form, giving 22050 or 11025.

\newpage
\section*{5. Topic-Wise Distribution}

\begin{longtable}{|l|c|l|}
\hline
\textbf{Topic} & \textbf{Number of Questions} & \textbf{Question Numbers} \\ \hline
Arithmetic (P/L, TSD, Mixtures, Boats, Averages, Time \& Work) & 6 & 1, 2, 6, 10, 12, 13 \\ \hline
Algebra (Logs, Quadratics, Functions, Inequalities, Series) & 6 & 3, 7, 9, 14, 15, (also 11 borderline) \\ \hline
Number System (Remainders, Divisibility) & 2 & 4, 8 \\ \hline
Geometry / Coordinate Geometry & 2 & 5, 11 \\ \hline
\end{longtable}

\section*{6. Quality Verification Summary}

\begin{longtable}{|c|c|c|c|c|}
\hline
\textbf{Question} & \textbf{Answer Verified} & \textbf{MCQ Options Verified} & \textbf{Ambiguity Check} & \textbf{Status} \\ \hline
1 & Yes & Yes & Clear & Ready \\ \hline
2 & Yes & Yes & Clear & Ready \\ \hline
3 & Yes & Yes & Clear & Ready \\ \hline
4 & Yes & Not Applicable & Clear & Ready \\ \hline
5 & Yes & Yes & Clear & Ready \\ \hline
6 & Yes & Not Applicable & Clear & Ready \\ \hline
7 & Yes & Not Applicable & Clear & Ready \\ \hline
8 & Yes & Yes & Clear & Ready \\ \hline
9 & Yes & Yes & Clear & Ready \\ \hline
10 & Yes & Not Applicable & Clear & Ready \\ \hline
11 & Yes & Yes & Clear & Ready \\ \hline
12 & Yes & Not Applicable & Clear & Ready \\ \hline
13 & Yes & Not Applicable & Clear & Ready \\ \hline
14 & Yes & Yes & Clear & Ready \\ \hline
15 & Yes & Not Applicable & Clear & Ready \\ \hline
\end{longtable}

\end{document}
